\documentclass{situationist-preprint}

\title{Colour Relationalism Defended}
\msid{MET-17-a}
\msfamily{MET-17}
\msdiscipline{Metaphysics}
\mstype{Principal manuscript}
\mskeywords{colour; relationalism; perceptual variation; eliminativism; physicalism; primitivism; response-dependence; secondary qualities}

\begin{document}
\maketitle

\begin{abstract}
Are colours real? Eliminativists hold that nothing is coloured and that colour experience is a systematic illusion; physicalists identify colours with perceiver-independent physical properties of surfaces; primitivists take them to be simple, non-physical, perceiver-independent qualities. We argue that colours are real, that they are relational properties of objects, and that they depend on perceivers without being illusory. The argument starts from perceptual variation. Normal perceivers, and one perceiver in different normal circumstances, experience the same surface as having incompatible colours, and nothing selects one of these experiences as the one that gets the colour right. If colours were perceiver-independent, at least one of the experiences would be illusory and there would be no fact about which. If colours are relational, both are veridical, each of the colour a surface has relative to a perceiver and a circumstance. Ordinary colour talk is context-sensitive, with the relevant perceivers and circumstances fixed by default to normal humans in daylight. Colours so understood are constructed, extrinsic properties: real, but not among the sparse universals. We reply to the objections from phenomenology, from disagreement, from proliferation and from the possibility of unknowable facts about true colours.
\end{abstract}

\section{The problem}

Locke distinguished the primary qualities of bodies, such as shape, size and motion, which resemble the ideas they produce in us, from the secondary qualities, such as colours, which are nothing in the objects but powers to produce sensations in us by means of their primary qualities \citep[II.viii]{Locke1690}. The distinction set the terms of a debate that is still open. Four positions are now prominent.

\emph{Eliminativism} holds that nothing has the colours we experience objects as having. \citet{Hardin1988} argued from colour science that nothing in the physical world corresponds to the structure of colour experience, and \citet{BoghossianVelleman1989} argued that colour experience represents colours as intrinsic, non-dispositional properties of surfaces, which nothing has, so that colour experience is in error.

\emph{Physicalism} identifies each colour with a physical property of surfaces, such as a type of surface spectral reflectance, which objects have independently of perceivers \citep{ByrneHilbert2003}. Physicalists must group physically heterogeneous reflectances under one colour, since many different reflectances look the same to us, and must say which of the many ways of grouping them is correct.

\emph{Primitivism} holds that colours are simple qualities, distinct from any physical property, that objects have independently of perceivers and that experience reveals as they are \citep{Campbell1993}.

\emph{Dispositionalism} and its generalization, \emph{relationalism}, hold that colours are constituted by relations between objects and perceivers. Locke's powers are dispositions to produce experiences; \citet{McGinn1983} and \citet{Johnston1992} developed response-dependent accounts; and \citet{Cohen2009} argued that colours are relational properties in general, of which dispositions are one kind.

We defend relationalism. Colours are real: ordinary colour attributions are, for the most part, true. They are relational: an object's colour is a relation between it, a type of perceiver and a type of circumstance. They are perceiver-dependent: which colour properties there are is fixed by the visual constitutions of possible perceivers. And they are not illusory: normal colour experience represents objects as having colours they have.

\section{Framework}

\begin{definition}[Relational colour]\label{d:relcol}
For a colour $C$, a type of perceiver $S$ and a type of circumstance $K$, an object $x$ is \emph{$C$ for $S$ in $K$} if and only if $x$ is disposed to look $C$ to perceivers of type $S$ in circumstances of type $K$. A \emph{relational colour property} is a property of this form.
\end{definition}

Types of perceiver are individuated by visual constitution: the photoreceptors, opponent processes and higher processing that determine how a perceiver's colour experiences covary with the light reaching the eye. Types of circumstance are individuated by illumination, surround and the state of adaptation. \enquote{Looks $C$} is to be read in its phenomenal sense: an object looks $C$ to a perceiver when her experience has the character that experiences of $C$ things have for her kind of perceiver.

\begin{principle}[Non-error]\label{p:nonerror}
Ordinary colour attributions, made by normal perceivers on the basis of perception in normal circumstances, are for the most part true.
\end{principle}

\begin{principle}[Variation]\label{p:variation}
For many objects $x$, there are pairs of normal perceivers, or pairs of normal circumstances, such that $x$ looks to have one determinate colour in the first and a different and incompatible determinate colour in the second.
\end{principle}

\begin{principle}[No privilege]\label{p:noprivilege}
In the cases of Principle~\ref{p:variation}, no fact about the object, the perceivers or the circumstances makes one of the experiences correct and the other incorrect.
\end{principle}

Principle~\ref{p:nonerror} is the datum the theory of colour must save, and the eliminativist's denial of it is the cost her theory must justify. Principle~\ref{p:variation} is established by colour science. Normal trichromats, who pass every standard test of colour vision, differ in the wavelength they select as unique green, the green that looks neither yellowish nor bluish, over a substantial range of the spectrum \citep{Hardin1988}. A chip that looks unique green to one such perceiver looks slightly bluish to another and slightly yellowish to a third. Variation across circumstances is more familiar: the same grey patch looks lighter against a dark surround than against a light one, and the same surface looks different under daylight and under incandescent light. Variation across species is more radical still, since animals with four types of cone discriminate surfaces that look identical to us.

Principle~\ref{p:noprivilege} is the substantive premise. In the unique green case, each perceiver is normal by every available criterion; the selection of one as the standard would be stipulation, and no physical fact about the chip favors one of them, since the chip's reflectance is the same whoever looks at it \citep{Cohen2009}. In the case of circumstances, no illumination is privileged by the nature of the surfaces it falls on: daylight varies across the day, and the visual system adapts to it. One might select the statistically typical perceiver and the most frequent circumstance, but statistical typicality is not a criterion of correctness, and the choice of a population and a range of circumstances over which to compute it is itself arbitrary.

\section{The argument}

\begin{lemma}[Variation forces error or relativity]\label{l:dilemma}
If each object has at most one determinate colour at a time, independently of perceivers and circumstances, then in every case of Principle~\ref{p:variation} at least one of the two experiences is illusory.
\end{lemma}

\begin{proof}
In such a case the object looks to have two incompatible determinate colours. If it has at most one determinate colour, it lacks at least one of them, and the experience presenting it as having that colour presents it as it is not.
\end{proof}

\begin{lemma}[No fact about which]\label{l:nofact}
If each object has at most one determinate colour, independently of perceivers and circumstances, then in the cases of Principle~\ref{p:variation} the object's colour is fixed by no fact about the object, the perceivers or the circumstances.
\end{lemma}

\begin{proof}
By Lemma~\ref{l:dilemma}, one of the experiences is in error. Any fact that fixed the object's colour would make the experience presenting that colour correct and the other incorrect, and by Principle~\ref{p:noprivilege} there is no such fact.
\end{proof}

\begin{theorem}[Relationalism]\label{t:main}
Colours are relational properties of objects: an object has a determinate colour only relative to a type of perceiver and a type of circumstance. In the cases of Principle~\ref{p:variation}, both experiences are veridical.
\end{theorem}

\begin{proof}
If each object had at most one determinate colour independently of perceivers and circumstances, its colour would, by Lemma~\ref{l:nofact}, be fixed by no fact about the object, the perceivers or the circumstances. Since nothing else could fix it, the object would have no determinate colour at all, and the colour attributions of perceivers in such cases would be false. Since such cases are pervasive, this would contradict Principle~\ref{p:nonerror}. So objects do not have their colours independently of perceivers and circumstances. If colours are relational in the sense of Definition~\ref{d:relcol}, the two experiences in a case of variation present the object as having two different relational properties, $C_1$ for $S_1$ in $K_1$ and $C_2$ for $S_2$ in $K_2$, which are compatible, and the object has both, since it is disposed to look $C_1$ to perceivers of the first type in the first circumstance and $C_2$ to perceivers of the second type in the second. Both experiences are veridical.
\end{proof}

The argument is a version of the argument from perceptual variation given by \citet{Cohen2009}. Its structure is that of an inference to the only view consistent with the three principles: eliminativism denies Principle~\ref{p:nonerror}; physicalism and primitivism must deny Principle~\ref{p:noprivilege}, since they hold that there is a fact about which perceiver gets the colour right; relationalism accepts all three.

\begin{proposition}[Perceiver-dependence]\label{p:dependence}
Which colour properties there are is fixed by the visual constitutions of possible perceivers. Had there been no possible perceivers with colour vision, nothing would have been coloured.
\end{proposition}

\begin{proof}
By Theorem~\ref{t:main}, every colour property is relative to a type of perceiver, and by Definition~\ref{d:relcol} a type of perceiver is individuated by its visual constitution. A colour relative to a type that no possible perceiver could instantiate would be a disposition to produce an experience that no possible being could have, and so would not be a colour property at all.
\end{proof}

Perceiver-dependence in this sense is not dependence on actual perceivers. An unseen flower in a world without animals is red for normal humans in daylight, because it is disposed to look red to such perceivers in such circumstances, whether or not any exist. What depends on perceivers is which properties are colours, not whether a given object, at a given time, is being perceived.

\begin{corollary}[Not illusory]\label{c:veridical}
Normal colour experience is not illusory. A perceiver's experience of an object as $C$ is veridical if the object is $C$ for her type of perceiver in her circumstances, and in normal cases it is.
\end{corollary}

\begin{proof}
An experience is illusory if it presents its object as having a property the object lacks. By Theorem~\ref{t:main}, the colour an experience presents is a colour relative to the perceiver's type and circumstances, and an object that looks $C$ to a normal perceiver of type $S$ in normal circumstances of type $K$ is, by the reliability of the visual system in such circumstances, disposed to look $C$ to perceivers of type $S$ in $K$. It is therefore $C$ for $S$ in $K$.
\end{proof}

\subsection{Ordinary colour talk}

When we say that a tomato is red, we do not mention perceivers or circumstances. On the relational view, ordinary colour predicates are context-sensitive: in a context of utterance, \enquote{red} expresses being red for the perceivers and circumstances salient in the context. By default these are normal human perceivers, as they actually are, in daylight. The indexing to actual human perceivers is rigid: under the counterfactual supposition that human visual systems were altered, the tomato would still be red in our sense, since it would still be disposed to look red to perceivers constituted as we actually are. This is the standard treatment of response-dependent predicates \citep{Johnston1992}, and it is the treatment AES-02 adopts for beauty, which it models on colour as understood here. Departures from the default are common and easily signalled: a lighting designer speaks of how a fabric will look under stage lights, and a biologist of the colours of flowers for bees.

\subsection{Colours as constructed properties}

The sparse theory of universals defended in MET-02 distinguishes the few genuine universals from the properties constructed from them. Relational colours belong to the second class. The property of being red for $S$ in $K$ is the property of having a surface whose genuine physical properties are such that, in circumstances of type $K$, it would produce in a perceiver of type $S$ an experience of a certain kind. Every element in this description is a genuine property or relation, or is constructed from such: surface reflectance from the physics of surfaces, the perceiver type from the physics and physiology of visual systems, and the disposition from the laws relating them. Colours are therefore constructed properties, in the sense of MET-02, and extrinsic ones, since whether an object has a given relational colour depends on how things other than the object, namely perceivers of the relevant type, are constituted. Their being constructed does not make them less real: they are instantiated, they figure in true attributions and in explanations of behavior, and construction is not elimination. Nor does their being dispositional, since dispositions are real properties of their bearers (MET-11). What their being constructed denies is only that they are among the fundamental joints of nature, which is the eliminativist's insight, rightly located.

\section{Objections and replies}

\begin{objection}[Phenomenology]
Colours do not look relational. A red surface looks red all over, intrinsically and categorically, not as disposed to produce experiences in some type of perceiver. If colour experience represents colours as monadic and intrinsic, then relationalism makes all colour experience inaccurate, and is a form of eliminativism after all \citep{BoghossianVelleman1989}.
\end{objection}

\begin{reply}
Colour experience presents colours as features of surfaces, and relationalism agrees that they are. What experience does not do is present them as non-relational. That a property is relational is not something experience reveals or denies; the claim that colours are relational is a theory of what is experienced, not a report of how it seems. Shape is a useful comparison: whether a coin looks round depends on the angle of view, and the property of looking elliptical from here is relational, yet no one concludes that experiences of perspectival shape are inaccurate. The eliminativist's argument requires the stronger premise that experience represents colours as being non-relational, and nothing in the phenomenology supports that premise rather than the weaker claim that experience does not represent colours as relational.
\end{reply}

\begin{objection}[True colours]
There may be a fact about which perceiver sees the true unique green, even if we cannot know it. Ignorance of the fact is no evidence of its absence, and the physicalist may hold that one perceiver is right and the others slightly wrong \citep{Tye2006}.
\end{objection}

\begin{reply}
The issue is not whether facts can be unknowable. MET-18 accepts facts of that kind: vague predicates have sharp boundaries that we cannot know. But in the vagueness case the boundary is fixed, though unknowably, by facts about the use of the predicate. In the case of unique green, the putative fact would have to select one of several perceivers whose use of \enquote{green} is equally competent and whose visual systems are equally normal, and no fact about use, about the chip or about the perceivers does so. A fact that supervenes on nothing is not an unknowable fact but no fact \citep{CohenHardinMcLaughlin2006}. The relationalist can and does accept the epistemicist's point at its proper place: for each type of perceiver and circumstance, the boundary of being green for that type in that circumstance may be sharp and unknowable.
\end{reply}

\begin{objection}[Disagreement]
Two perceivers who say \enquote{unique green} and \enquote{slightly bluish} of the same chip disagree. On relationalism, each attributes a different relational property, and there is no disagreement. Relationalism misdescribes the case.
\end{objection}

\begin{reply}
When the perceivers talk to each other, their context fixes a common index, typically normal human perceivers in the circumstances at hand, and under that index at most one of them is right. When they discover that they differ in visual constitution, the disagreement dissolves into an agreement that the chip looks one way to one and another way to the other, and this is what happens in practice: subjects in colour-matching experiments who learn that their settings differ do not continue to dispute which of them is correct. Relationalism predicts both phases of the exchange, the apparent disagreement and its resolution, while the views on which one perceiver must be right predict a disagreement that persists without any means of resolution.
\end{reply}

\begin{objection}[Proliferation]
Relationalism multiplies colours without limit. Every object has a different colour for every type of perceiver and circumstance, including perceivers of no actual species. That is an extravagant ontology.
\end{objection}

\begin{reply}
The colours are many, but they are constructed, and constructed properties cost nothing beyond the genuine properties from which they are constructed (MET-02). The ontology of the relationalist contains the surfaces, the perceivers and the laws, all of which every theory of colour accepts; the relational colours are what those entities determine. A theory of mass that recognizes a determinate mass for every real number of kilograms is not extravagant either.
\end{reply}

\begin{objection}[Inverted spectrum]
If two perceivers could be functionally identical while one's experiences of red are like the other's experiences of green, then relationalism, which individuates perceiver types by visual constitution, cannot say which colour each sees.
\end{objection}

\begin{reply}
The case concerns the relation between colour properties and the phenomenal character of experience, rather than the reality of colours, and it is treated in MET-17-b. For the present argument, it suffices that Definition~\ref{d:relcol} individuates colours by how objects look, in the phenomenal sense, to perceivers of a type, and that inverted perceivers, if possible, would belong to different types.
\end{reply}

\section{What the result does not establish}

The argument does not establish an account of the phenomenal character of colour experience, or of how experience represents colours. It does not establish that the other secondary qualities, sounds, tastes and smells, are relational, although the argument from variation applies to them with obvious modifications. It does not imply that colour science is mistaken in classifying surfaces by reflectance: reflectance types are the physical bases of relational colours, and they are grouped into colours by their effects on perceivers. And it does not imply that any property that depends on responses is relational in the same way; AES-02 argues that beauty is, but the argument there rests on its own principles.

\section{Conclusion}

Colours are real. Perceptual variation shows that, if each object had a single perceiver-independent colour, normal colour experience would be in error in pervasive cases with no fact to decide where the error lies. Relationalism avoids the dilemma: objects have colours relative to types of perceiver and circumstance, ordinary colour talk fixes the relevant type by context, and the experiences of normal perceivers in normal circumstances are veridical. Colours so understood depend on perceivers without being illusory, and they are constructed, extrinsic properties: real, though not among the fundamental constituents of the world.

\begin{thebibliography}{99}\small
\bibitem[Boghossian and Velleman(1989)]{BoghossianVelleman1989} Boghossian, P. A., and Velleman, J. D. (1989). Colour as a secondary quality. \emph{Mind}, 98, 81--103.
\bibitem[Byrne and Hilbert(2003)]{ByrneHilbert2003} Byrne, A., and Hilbert, D. R. (2003). Color realism and color science. \emph{Behavioral and Brain Sciences}, 26, 3--21.
\bibitem[Campbell(1993)]{Campbell1993} Campbell, J. (1993). A simple view of colour. In J. Haldane and C. Wright (eds.), \emph{Reality, Representation, and Projection}. Oxford University Press.
\bibitem[Cohen(2009)]{Cohen2009} Cohen, J. (2009). \emph{The Red and the Real: An Essay on Color Ontology}. Oxford University Press.
\bibitem[Cohen et~al.(2006)]{CohenHardinMcLaughlin2006} Cohen, J., Hardin, C. L., and McLaughlin, B. P. (2006). True colours. \emph{Analysis}, 66, 335--340.
\bibitem[Hardin(1988)]{Hardin1988} Hardin, C. L. (1988). \emph{Color for Philosophers: Unweaving the Rainbow}. Hackett.
\bibitem[Johnston(1992)]{Johnston1992} Johnston, M. (1992). How to speak of the colors. \emph{Philosophical Studies}, 68, 221--263.
\bibitem[Locke(1690)]{Locke1690} Locke, J. (1690). \emph{An Essay Concerning Human Understanding}. London.
\bibitem[McGinn(1983)]{McGinn1983} McGinn, C. (1983). \emph{The Subjective View: Secondary Qualities and Indexical Thoughts}. Oxford University Press.
\bibitem[Tye(2006)]{Tye2006} Tye, M. (2006). The puzzle of true blue. \emph{Analysis}, 66, 173--178.
\end{thebibliography}

\end{document}
